\section{Introduction}

Ce document décrit la \textbf{procédure de configuration}, sous l'atelier
logiciel Control Expert, d'un coupleur Ethernet/IP \textbf{BMX~NOC~0401.2}
équipant l'automate Schneider M340 de la cellule robotique. Cette
procédure est établie à partir du support \emph{RIF\_GT\_TP1\_2024}, qui
l'illustre sur les deux réseaux effectivement présents sur l'installation :

\begin{itemize}
    \item le réseau \textbf{NOC\_ETHIP\_CONVOYEUR}, qui scrute le
    contrôleur PFC200 assurant la commande du poste guichet ;
    \item le réseau \textbf{NOC\_ETHIP\_ROBOT}, qui scrute la baie CS9 du
    robot et le contrôleur M221.
\end{itemize}

\begin{UPSTIinfor}{Objet du document}
La procédure décrite ici est \textbf{générique} : elle s'applique à la
configuration de tout coupleur BMX~NOC~0401.2 destiné à scruter un ou
plusieurs équipements Ethernet/IP. Elle est illustrée, à chaque étape, par
son application au réseau \texttt{NOC\_ETHIP\_CONVOYEUR} (scrutation du
PFC200), qui constitue l'exemple le plus complet et le plus simple à
suivre.
\end{UPSTIinfor}

\subsection{Étapes de la procédure}

La configuration d'un coupleur NOC et de son réseau associé s'effectue en
cinq étapes, à réaliser dans l'ordre :

\begin{enumerate}
    \item \textbf{Importation des fichiers EDS} décrivant les équipements
    à scruter (préalable, une seule fois par équipement) ;
    \item \textbf{Configuration matérielle} du coupleur dans le rack de
    l'automate (emplacement, nom du DTM, protocole) ;
    \item \textbf{Configuration réseau} du coupleur (adressage IP) ;
    \item \textbf{Ajout et configuration d'un équipement} sur le réseau
    ainsi créé (DTM de l'équipement, paramètres d'échange, adressage,
    nommage des données) ;
    \item \textbf{Vérification} de la configuration obtenue (compilation,
    éditeur de données, contrôle de la cartographie mémoire) avant tout
    téléchargement dans l'automate.
\end{enumerate}

Chacune de ces étapes est détaillée dans les sections suivantes.
